\chapter{सम्मिश्र संख्याएँ}\label{ch:g12:complex}

\omterm{def:g10:algebra:equation}{समीकरण} $x^2 = -1$ का कोई वास्तविक हल नहीं है। $\R$ को केवल एक नई संख्या $\iu$ से बढ़ा
देने पर, जिसका वर्ग $-1$ है, \omterm{def:g12:complex:def}{सम्मिश्र संख्याओं} का क्षेत्र $\C$ बन जाता है,
जिसमें \emph{हर} बहुपद \omterm{def:g10:algebra:equation}{समीकरण} के हल होते हैं --- और जिसका अंकगणित समतल की
ज्यामिति समेटे रहता है: स्थानांतरण, घूर्णन और समरूपण योग तथा गुणन बन जाते
हैं।

\section{बीजगणितीय रूप}

\begin{definition}[सम्मिश्र संख्याएँ]\label{def:g12:complex:def}
\emph{सम्मिश्र संख्याओं}\index{सम्मिश्र संख्या} का समुच्चय
\[
\C = \{a + \iu b : a, b \in \R\},
\]
है, जहाँ $\iu$ एक ऐसा संकेत है जिस पर केवल नियम $\iu^2 = -1$ लागू है; योग और गुणन $\iu$
में वास्तविक बहुपदों की तरह किए जाते हैं और $\iu^2$ को $-1$ तक घटा दिया जाता है।
\omterm{def:g10:numbers:sets}{वास्तविक संख्याएँ} $a$ और $b$ $z = a + \iu b$ के \emph{वास्तविक भाग}\index{वास्तविक भाग}
$\Rea(z)$ और \emph{काल्पनिक भाग}\index{काल्पनिक भाग} $\Ima(z)$ हैं; यह व्यंजक $z$ का
\emph{बीजगणितीय रूप}\index{बीजगणितीय रूप} है, और वह अद्वितीय है: $a + \iu b = a' + \iu b'$ यदि और
केवल यदि $a = a'$ और $b = b'$ हों।
\end{definition}

\begin{example}
$(2 + 3\iu)(1 - \iu) = 2 - 2\iu + 3\iu - 3\iu^2 = 5 + \iu$।
\end{example}

\begin{definition}[संयुग्मी]\label{def:g12:complex:conjugate}
$z = a + \iu b$ का \emph{संयुग्मी}\index{संयुग्मी} $\conj{z} = a - \iu b$ है।
\end{definition}

\begin{proposition}[संयुग्मन के गुण]\label{prop:g12:complex:conjugate}
सभी $z, w \in \C$ के लिए:
\[
\conj{z + w} = \conj z + \conj w, \quad
\conj{zw} = \conj z\,\conj w, \quad
\conj{\conj z} = z, \quad
z + \conj z = 2\Rea(z), \quad
z - \conj z = 2\iu\,\Ima(z),
\]
और $z = a + \iu b$ के लिए $z\conj z = a^2 + b^2 \geq 0$। इसके अतिरिक्त $z \in \R$ यदि और केवल यदि $z = \conj z$ हो।
\end{proposition}

\begin{proof}
सभी \omterm{def:g12:complex:def}{बीजगणितीय रूपों} से सीधी गणनाएँ हैं; जैसे $z\conj z = (a + \iu b)(a - \iu b) = a^2 - (\iu b)^2 = a^2 + b^2$।
\end{proof}

\begin{method}[सम्मिश्र संख्याओं का भाग]\label{met:g12:complex:division}
किसी भागफल को \omterm{def:g12:complex:def}{बीजगणितीय रूप} में लिखने के लिए अंश और हर दोनों को हर के \omterm{def:g12:complex:conjugate}{संयुग्मी}
से गुणा कीजिए:
\[
\frac{1}{2 + \iu} = \frac{2 - \iu}{(2+\iu)(2-\iu)} = \frac{2 - \iu}{5}
= \frac25 - \frac15\iu .
\]
इसलिए हर शून्येतर \omterm{def:g12:complex:def}{सम्मिश्र संख्या} का प्रतिलोम होता है: $\C$ एक \emph{क्षेत्र}
है।
\end{method}

\section{सम्मिश्र समतल, मापांक और कोणांक}

\begin{definition}[बिंदुसंख्या, मापांक]\label{def:g12:complex:modulus}
\omterm{def:g10:coordgeom:system}{लंबकोणिक} \omterm{def:g10:coordgeom:system}{निर्देशांकों} वाले समतल में बिंदु $M(a, b)$ $z = a + \iu b$ का \emph{\omterm{def:g10:functions:function}{प्रतिबिंब}} है,
और $z$ $M$ की \emph{बिंदुसंख्या}\index{बिंदुसंख्या} है। $z$ का
\emph{मापांक}\index{मापांक} $M$ की मूल बिंदु से दूरी है:
\[
\abs{z} = \sqrt{a^2 + b^2} = \sqrt{z\conj z}.
\]
\end{definition}

\begin{omfigure}
\begin{tikzpicture}[scale=1.5]
  \draw[->] (-0.7,0) -- (2.9,0) node[right] {$x$};
  \draw[->] (0,-1.9) -- (0,2.1) node[above] {$y$};
  \draw[dashed] (2,1.5) -- (2,0) node[below right] {$a$};
  \draw[dashed] (2,1.5) -- (0,1.5) node[left] {$b$};
  \draw[dashed] (2,-1.5) -- (2,0);
  \draw[-Stealth,omDef,thick] (0,0) -- (2,1.5)
    node[midway, above left, font=\small] {$\abs z$};
  \draw[-Stealth,omThm,thick] (0,0) -- (2,-1.5);
  \fill (2,1.5) circle (1pt) node[above right] {$M(z = a + \iu b)$};
  \fill (2,-1.5) circle (1pt) node[below right] {$\conj z = a - \iu b$};
\end{tikzpicture}

{\small \omterm{def:g12:complex:modulus}{बिंदुसंख्या} $z$ वाला बिंदु $M(a,b)$: \omterm{def:g12:complex:modulus}{मापांक} $\abs z$ लंबाई $OM$ है, और
\omterm{def:g12:complex:conjugate}{संयुग्मी} $\conj z$ वास्तविक अक्ष में $z$ का परावर्तन है।}
\end{omfigure}

\begin{proposition}[मापांक के गुण]\label{prop:g12:complex:modulus}
$z, w \in \C$ के लिए:
\[
\abs{zw} = \abs z\,\abs w, \qquad
\abs{\tfrac zw} = \tfrac{\abs z}{\abs w}\ (w \neq 0), \qquad
\abs{\conj z} = \abs z, \qquad
\abs{z + w} \leq \abs{z} + \abs{w} \ \text{(त्रिभुज असमिका)}.
\]
इसके अतिरिक्त $\abs{z_B - z_A}$ बिंदुसंख्याओं $z_A$ और $z_B$ वाले बिंदुओं के बीच की दूरी है।
\end{proposition}

\begin{proof}
$\abs{zw}^2 = zw\,\conj{zw} = z\conj z\, w \conj w = \abs z^2 \abs w^2$
से पहला मिलता है; दूसरा और तीसरा इसी जैसे हैं। त्रिभुज असमिका के लिए:
\[
\abs{z+w}^2 = (z+w)(\conj z + \conj w)
= \abs z^2 + \abs w^2 + 2\Rea(z \conj w)
\leq \abs z^2 + \abs w^2 + 2\abs{z\conj w}
= \bigl(\abs z + \abs w\bigr)^2,
\]
जहाँ $\Rea(u) \leq \abs u$ का उपयोग हुआ है। दूरी वाला कथन $z_B - z_A$ के \omterm{def:g10:coordgeom:system}{निर्देशांकों} पर लगाया गया
पाइथागोरस सूत्र है।
\end{proof}

\begin{definition}[कोणांक, त्रिकोणमितीय और चरघातांकी रूप]\label{def:g12:complex:argument}
मान लीजिए $z \neq 0$ है। $z$ का \emph{कोणांक}\index{कोणांक}, जिसे $\Arg(z)$ लिखते हैं,
धनात्मक वास्तविक अक्ष से किरण $OM$ तक के कोण का कोई भी माप $\theta$ है; वह $2k\pi$
तक ही परिभाषित होता है। $r = \abs z$ लिखने पर
\[
z = r(\cos\theta + \iu\sin\theta)
\qquad\text{(त्रिकोणमितीय रूप)}.
\]
संकेतन\index{चरघातांकी रूप}
\[
\eu^{\iu\theta} = \cos\theta + \iu\sin\theta,
\]
प्रस्तुत करने पर यह \emph{चरघातांकी रूप} $z = r\,\eu^{\iu\theta}$ बन जाता है।
\end{definition}

\begin{omfigure}
\begin{tikzpicture}[scale=1.5]
  \draw[->] (-0.7,0) -- (2.4,0) node[right] {$x$};
  \draw[->] (0,-0.5) -- (0,1.9) node[above] {$y$};
  \draw[gray, dashed] (1.9,0) arc (0:75:1.9);
  \draw[-Stealth,omDef,thick] (0,0) -- (40:1.9)
    node[pos=0.62, below right=1pt, font=\small] {$r = \abs z$};
  \fill (40:1.9) circle (1pt) node[above right] {$z = r\,\eu^{\iu\theta}$};
  \draw[->,omThm,thick] (0.45,0) arc (0:40:0.45);
  \node[omThm] at (20:0.65) {$\theta$};
\end{tikzpicture}

{\small \omterm{def:g12:complex:argument}{चरघातांकी रूप}: $z$ उसकी मूल बिंदु से दूरी $r$ और धनात्मक वास्तविक
अक्ष से कोण $\theta$ द्वारा तय हो जाता है।}
\end{omfigure}

\begin{theorem}\label{thm:g12:complex:expform}
सभी $\theta, \theta' \in \R$ के लिए:
\[
\eu^{\iu\theta}\,\eu^{\iu\theta'} = \eu^{\iu(\theta + \theta')},
\qquad
\conj{\eu^{\iu\theta}} = \eu^{-\iu\theta},
\qquad
\abs{\eu^{\iu\theta}} = 1 .
\]
परिणामस्वरूप, शून्येतर $z = r\eu^{\iu\theta}$ और $z' = r'\eu^{\iu\theta'}$ के लिए:
\[
zz' = rr'\,\eu^{\iu(\theta+\theta')},
\qquad
\frac{z}{z'} = \frac{r}{r'}\,\eu^{\iu(\theta - \theta')} :
\]
\emph{\omterm{def:g12:complex:modulus}{मापांक} गुणा होते हैं, \omterm{def:g12:complex:argument}{कोणांक} जुड़ जाते हैं।}
\end{theorem}

\begin{proof}
पहली सर्वसमिका योग सूत्रों (\cref{prop:g12:trigo:addition}) का युग्म \emph{ही} है:
\[
(\cos\theta + \iu\sin\theta)(\cos\theta' + \iu\sin\theta')
= (\cos\theta\cos\theta' - \sin\theta\sin\theta')
+ \iu(\sin\theta\cos\theta' + \cos\theta\sin\theta').
\]
शेष सीधी गणना से निकलता है।
\end{proof}

\begin{corollary}[दे म्वाव्र का सूत्र]\label{cor:g12:complex:demoivre}
\index{दे म्वाव्र का सूत्र}
$\theta \in \R$ और $n \in \Z$ के लिए: $\bigl(\cos\theta + \iu\sin\theta\bigr)^n = \cos n\theta + \iu \sin n\theta$।
\end{corollary}

\begin{proof}
\cref{thm:g12:complex:expform} का उपयोग करके $n \geq 0$ पर आगमन; प्रतिलोम लेकर $n < 0$ तक बढ़ाइए।
\end{proof}

\begin{method}[रूपों के बीच आना-जाना]\label{met:g12:complex:forms}
बीजगणितीय से चरघातांकी: $r = \abs z$ निकालिए, फिर ऐसा $\theta$ ज्ञात कीजिए कि $\cos\theta = \frac ar$, $\sin\theta = \frac br$
हों (प्रतिलोम \omterm{def:g12:trigo:cossin}{कोज्या} का हवाला देने से पहले सही चतुर्थांश पहचान लीजिए)।
चरघातांकी से बीजगणितीय: $r\cos\theta + \iu\, r\sin\theta$ का प्रसार कीजिए। योगों के लिए \omterm{def:g12:complex:def}{बीजगणितीय रूप}, और
गुणनफलों, घातों तथा भागफलों के लिए \omterm{def:g12:complex:argument}{चरघातांकी रूप} काम में लाइए।
\end{method}

\begin{example}
$z = 1 + \iu$: $\abs z = \sqrt2$ और $\cos\theta = \sin\theta =
\frac{1}{\sqrt2}$, इसलिए $\theta = \frac\pi4$ तथा $z = \sqrt2\,\eu^{\iu\pi/4}$. अतः $z^{20} = 2^{10}\,\eu^{5\iu\pi} = -1024$।
\end{example}

\section{द्विघात समीकरण और इकाई के मूल}

\begin{theorem}[वास्तविक गुणांकों वाले द्विघात समीकरण]\label{thm:g12:complex:quadratic}
मान लीजिए $a, b, c \in \R$, $a \neq 0$, और $\Delta = b^2 - 4ac$ हैं। \omterm{def:g10:algebra:equation}{समीकरण} $az^2 + bz + c = 0$ के $\C$ में ये हल हैं:
\[
\begin{aligned}
&\Delta > 0:\ z = \frac{-b \pm \sqrt{\Delta}}{2a} \ (\text{दो वास्तविक}); \qquad
\Delta = 0:\ z = \frac{-b}{2a} \ (\text{द्विगुण}); \\
&\Delta < 0:\ z = \frac{-b \pm \iu\sqrt{-\Delta}}{2a} \ (\text{दो संयुग्मी}).
\end{aligned}
\]
\end{theorem}

\begin{proof}
वर्ग पूरा कीजिए: $az^2 + bz + c = a\left(\left(z + \frac{b}{2a}\right)^2 -
\frac{\Delta}{4a^2}\right)$। यदि $\Delta < 0$ हो, तो $\frac{\Delta}{4a^2} = \left(\iu\frac{\sqrt{-\Delta}}{2a}\right)^2$, और वर्गों का अंतर सामान्य
तरीक़े से गुणनखंडित हो जाता है।
\end{proof}

\begin{definition}[इकाई के मूल]\label{def:g12:complex:rootsofunity}
$n \geq 1$ के लिए \emph{इकाई के $n$-वें मूल}\index{इकाई के मूल} $z^n = 1$ के हल हैं।
\end{definition}

\begin{theorem}\label{thm:g12:complex:rootsofunity}
\omterm{def:g10:algebra:equation}{समीकरण} $z^n = 1$ के ठीक $n$ हल हैं:
\[
\omega_k = \eu^{2\iu k\pi/n}, \qquad k = 0, 1, \dots, n-1 .
\]
उनके \omterm{def:g10:functions:function}{प्रतिबिंब} इकाई वृत्त में अंतर्लिखित सम $n$-भुज के शीर्ष हैं, और $n \geq 2$
के लिए उनका योग $0$ है।
\end{theorem}

\begin{proof}
हर $\omega_k$ $\omega_k^n = \eu^{2\iu k\pi} = 1$ को संतुष्ट करता है, और वे जोड़े-जोड़े भिन्न हैं क्योंकि उनके
\omterm{def:g12:complex:argument}{कोणांक} $\intco{0}{2\pi}$ में पड़ते हैं। इसके उलट, यदि $z^n = 1$ हो, तो $\abs z^n = 1$, इसलिए $\abs z = 1$
(धनात्मक वास्तविक) और $n\theta \equiv 0 \pmod{2\pi}$ वाला $z = \eu^{\iu\theta}$, \emph{अर्थात्} $\theta = \frac{2k\pi}{n}$; $k$ को $n$
के सापेक्ष घटाने पर वह $\omega_k$ में से किसी एक पर आ जाता है। शीर्ष कोण $\frac{2\pi}{n}$ से
बराबर-बराबर बँटे हैं: यानी सम $n$-भुज। अंत में, $\omega = \omega_1$ के साथ $\omega \neq 1$ और
\omterm{def:g12:seq:arith-geom}{गुणोत्तर} योग देता है
\[
\sum_{k=0}^{n-1}\omega_k = \sum_{k=0}^{n-1}\omega^k
= \frac{\omega^n - 1}{\omega - 1} = 0 . \qedhere
\]
\end{proof}

\begin{omfigure}
\begin{tikzpicture}[scale=1.7]
  \draw[->] (-1.3,0) -- (1.45,0) node[right] {$x$};
  \draw[->] (0,-1.3) -- (0,1.4) node[above] {$y$};
  \draw[gray] (0,0) circle (1);
  \draw[omDef] (0:1) -- (72:1) -- (144:1) -- (216:1) -- (288:1) -- cycle;
  \draw[->,omThm,thick] (0.3,0) arc (0:72:0.3);
  \node[omThm, font=\small] at (36:0.52) {$\frac{2\pi}{5}$};
  \fill (0:1) circle (1pt) node[below right] {$1$};
  \fill (72:1) circle (1pt) node[above right] {$\omega$};
  \fill (144:1) circle (1pt) node[above left] {$\omega^2$};
  \fill (216:1) circle (1pt) node[below left] {$\omega^3$};
  \fill (288:1) circle (1pt) node[below right] {$\omega^4$};
\end{tikzpicture}

{\small इकाई के पाँचवें \omterm{def:g11:quad:discriminant}{मूल}, $\omega = \eu^{2\iu\pi/5}$: इकाई वृत्त में अंतर्लिखित सम पंचभुज के
शीर्ष (\cref{exo:g12:complex:10})।}
\end{omfigure}

\begin{example}
इकाई के घनमूल $1$, $j = \eu^{2\iu\pi/3} =
-\frac12 + \iu\frac{\sqrt3}{2}$ और $\conj j = j^2$ हैं, और वे $1 + j + j^2 = 0$ को संतुष्ट करते हैं।
\end{example}

\section{सम्मिश्र संख्याएँ और समतल ज्यामिति}

\begin{proposition}[ज्यामितीय शब्दकोश]\label{prop:g12:complex:geometry}
मान लीजिए $A, B, C, D$ बिंदुसंख्याओं $a, b, c, d$ ($a \neq b$, $c \neq d$) वाले बिंदु हैं।
\begin{enumerate}
  \item \omterm{def:g12:complex:modulus}{बिंदुसंख्या} $t$ वाले \omterm{def:g11:vect:vector}{सदिश} से स्थानांतरण $z \mapsto z + t$ है।
  \item केंद्र $\omega$ (\omterm{def:g12:complex:modulus}{बिंदुसंख्या}) और कोण $\theta$ वाला घूर्णन $z \mapsto \omega + \eu^{\iu\theta}(z - \omega)$ है।
  \item केंद्र $\omega$ और अनुपात $k \in \R^*$ वाला समरूपण $z \mapsto \omega + k(z - \omega)$ है।
  \item $\dfrac{d - c}{b - a}$ का \omterm{def:g12:complex:modulus}{मापांक} $\dfrac{CD}{AB}$ है और \omterm{def:g12:complex:argument}{कोणांक} \omterm{def:g10:vectors:vector}{सदिशों} $\vect{AB}$ तथा $\vect{CD}$ के बीच का
        कोण। विशेष रूप से, $AB \perp CD$ यदि और केवल यदि $\frac{d-c}{b-a}$ विशुद्ध काल्पनिक हो,
        और $A, B, C$ संरेख ($C \ne A$, $C\ne B$) हैं यदि और केवल यदि $\frac{c-a}{b-a} \in \R$ हो।
\end{enumerate}
\end{proposition}

\begin{proof}
(1) और (3) \omterm{def:g10:coordgeom:system}{निर्देशांकों} वाले सूत्रों को ही दोहराते हैं। (2) के लिए: प्रतिचित्रण
$w \mapsto \eu^{\iu\theta} w$ \omterm{def:g12:complex:modulus}{मापांकों} को $1$ से गुणा करता है और \omterm{def:g12:complex:argument}{कोणांकों} में $\theta$ जोड़ देता है:
यानी वह मूल बिंदु के परितः कोण $\theta$ का घूर्णन है; $\omega$ को $0$ पर भेजने
वाले स्थानांतरण से संयुग्मन करने पर सामान्य केंद्र मिल जाता है। (4) के लिए:
भागफलों के \omterm{def:g12:complex:modulus}{मापांक} और \omterm{def:g12:complex:argument}{कोणांक} घट जाते हैं (\cref{thm:g12:complex:expform})।
\end{proof}

\section{अभ्यास}

\begin{exercise}[$\star$]\label{exo:g12:complex:1}
\omterm{def:g12:complex:def}{बीजगणितीय रूप} में लिखिए:
$(3 - 2\iu)(1 + 4\iu)$, \quad $\dfrac{2 + \iu}{1 - \iu}$, \quad $\iu^{2027}$।
\end{exercise}

\begin{exercise}[$\star$]\label{exo:g12:complex:2}
$\C$ में हल कीजिए: (a) $z^2 - 4z + 13 = 0$; (b) $z^2 + z + 1 = 0$। हर स्थिति में जाँचिए कि दोनों हल
\omterm{def:g12:complex:conjugate}{संयुग्मी} हैं और उनका गुणनफल निकालिए।
\end{exercise}

\begin{exercise}[$\star$]\label{exo:g12:complex:3}
\omterm{def:g12:complex:argument}{चरघातांकी रूप} में लिखिए: $z_1 = -1 + \iu$, $z_2 = \sqrt3 - \iu$, और $\dfrac{z_1}{z_2}$ को पहले चरघातांकी फिर
\omterm{def:g12:complex:def}{बीजगणितीय रूप} में निकालिए। इससे $\cos\dfrac{11\pi}{12}$ का यथार्थ मान निकालिए।
\end{exercise}

\begin{exercise}[$\star\star$]\label{exo:g12:complex:4}
\omterm{def:g12:complex:modulus}{बिंदुसंख्या} $z$ वाले बिंदुओं $M$ के उन समुच्चयों का ज्यामितीय वर्णन कीजिए
जिनके लिए:
(a) $\abs{z - 2} = \abs{z + \iu}$;\quad
(b) $\abs{z - 1 - \iu} = 3$;\quad
(c) $\dfrac{z - 1}{z + 1}$ विशुद्ध काल्पनिक हो।
\end{exercise}

\begin{exercise}[$\star\star$]\label{exo:g12:complex:5}
दे म्वाव्र के सूत्र और द्विपद प्रमेय का उपयोग करके $\cos 3\theta$ को $\cos\theta$ के बहुपद के
रूप में, और $\sin 3\theta$ को $\sin\theta$ के पदों में व्यक्त कीजिए।
\end{exercise}

\begin{exercise}[$\star\star$]\label{exo:g12:complex:6}
\emph{(रैखिकीकरण.)} $\cos\theta = \frac{\eu^{\iu\theta} + \eu^{-\iu\theta}}{2}$ का उपयोग करके $\cos^3\theta$ का रैखिकीकरण कीजिए (उसे $\cos k\theta$
के संयोजन के रूप में लिखिए), और $\displaystyle\int_0^{\pi/2}\cos^3\theta\,\dd\theta$ निकालिए।
\end{exercise}

\begin{exercise}[$\star\star$]\label{exo:g12:complex:7}
मान लीजिए $A$ और $B$ की बिंदुसंख्याएँ $a = 1 + \iu$ और $b = 3 + 2\iu$ हैं। वे दोनों बिंदु
$C$ ज्ञात कीजिए जो त्रिभुज $ABC$ को समबाहु बनाते हैं, और उन्हें केंद्र $A$
तथा कोणों $\pm\frac{\pi}{3}$ वाले घूर्णनों के अधीन $B$ के \omterm{def:g10:functions:function}{प्रतिबिंब} के रूप में लीजिए।
\end{exercise}

\begin{exercise}[$\star\star\star$]\label{exo:g12:complex:8}
$\C$ में $z^4 = -4$ हल कीजिए और हल आलेखित कीजिए। $z^4 + 4$ का वास्तविक गुणांकों वाले
दो द्विघात बहुपदों में \omterm{def:g10:algebra:expand}{गुणनखंडन} कीजिए।
\end{exercise}

\begin{exercise}[$\star\star\star$]\label{exo:g12:complex:9}
$\theta \in \intoo{0}{2\pi}$ और $n \in \N$ के लिए \omterm{def:g12:seq:arith-geom}{गुणोत्तर} श्रेढ़ी $\sum_{k=0}^n \eu^{\iu k\theta}$ का योग निकालकर और \omterm{def:g12:complex:def}{वास्तविक भाग}
लेकर
\[
S = 1 + \cos\theta + \cos 2\theta + \dots + \cos n\theta
\]
निकालिए। (उत्तर: $S = \dfrac{\sin\frac{(n+1)\theta}{2}}{\sin\frac{\theta}{2}}
\cos\frac{n\theta}{2}$।)
\end{exercise}

\begin{exercise}[$\star\star\star$]\label{exo:g12:complex:10}
मान लीजिए $\omega = \eu^{2\iu\pi/5}$ है।
\begin{enumerate}
  \item पुष्ट कीजिए कि $1 + \omega + \omega^2 + \omega^3 + \omega^4 = 0$।
  \item मान लीजिए $u = \omega + \omega^4$ और $v = \omega^2 + \omega^3$ हैं। दिखाइए कि $u + v = -1$ और $uv = -1$, और $u = \frac{-1+\sqrt5}{2}$
        निकालिए।
  \item निष्कर्ष निकालिए कि $\cos\dfrac{2\pi}{5} = \dfrac{\sqrt5 - 1}{4}$।
\end{enumerate}
\end{exercise}

\section{समस्या: इकाई के दर्पण में पंचभुज}

\begin{problem}[{सप्ताहांत समस्या --- इकाई के पाँचवें मूल $\cos 72^\circ$ का यथार्थ मान
निकाल देते हैं, स्वर्ण अनुपात परिणाम पर हस्ताक्षर करता है, और गुणा वस्तुतः
घूर्णन निकलता है}]
\label{pb:g12:complex:1}
यूक्लिड सम पंचभुज की रचना कर सकते थे, पर वह पटरी और परकार के आगे \emph{क्यों}
झुक जाता है --- जबकि मामूली $20^\circ$ का कोण नहीं झुकता --- यह कारण दो हज़ार
वर्षों तक छिपा रहा। वह इसी अध्याय में छिपा है: इकाई के पाँचों पाँचवें \omterm{def:g11:quad:discriminant}{मूल}
(\cref{thm:g12:complex:rootsofunity}) का योग शून्य है, और बीजगणित की वही एक पंक्ति $\cos 72^\circ$ को एक द्विघात
\omterm{def:g10:algebra:equation}{समीकरण} में से निचोड़ लेती है --- केवल वर्गमूल, अतः रचनीय, और उस पर हस्ताक्षर
भी, ज़ाहिर है, स्वर्ण अनुपात के। इस रत्न के चारों ओर: प्रवाह, त्रिकोणमिति तक
दे म्वाव्र का छोटा रास्ता, और गुणा का ज्यामिति के रूप में पर्दाफ़ाश।

\medskip
\noindent\textbf{भाग I --- प्रवाह।}
\begin{enumerate}
  \item $(2 + \iu)(3 - \iu)$, \ $\dfrac{1 + \iu}{1 - \iu}$ (\cref{met:g12:complex:division}), और $\abs{3 + 4\iu}$ निकालिए।
  \item \omterm{def:g12:complex:argument}{चरघातांकी रूप} में लिखिए: $1 + \iu$; \ $-2$; \ $1 - \iu\sqrt3$।
  \item \omterm{def:g12:complex:argument}{चरघातांकी रूप} का उपयोग करके $(1 + \iu)^8$ निकालिए।
  \item $z^2 = \iu$ हल कीजिए, फिर $z^2 - 2z + 5 = 0$ (\cref{thm:g12:complex:quadratic})।
  \item ज्यामितीय शब्दकोश (\cref{prop:g12:complex:geometry}): प्रतिचित्रण $z \mapsto \iu z$ और समुच्चय $\abs{z - (1 + \iu)} = 2$ का
        वर्णन कीजिए।
\end{enumerate}

\medskip
\noindent\textbf{भाग II --- पाँच \omterm{def:g11:quad:discriminant}{मूल}, एक पंचभुज।}
मान लीजिए $\omega = \eu^{2\iu\pi/5}$ है और $1, \omega, \omega^2, \omega^3, \omega^4$ पर विचार कीजिए।
\begin{enumerate}[resume]
  \item जाँचिए कि ये ठीक $z^5 = 1$ के हल हैं, और समतल में वे जो आकृति खींचते हैं
        उसका वर्णन कीजिए।
  \item सिद्ध कीजिए कि उनका योग शून्य है। (\omterm{def:g12:seq:arith-geom}{गुणोत्तर} श्रेढ़ी --- या योग को
        $1 - \omega$ से गुणा कीजिए।)
  \item \omterm{def:g12:complex:def}{वास्तविक भाग} लेकर निकालिए:
        \[
        1 + 2\cos\frac{2\pi}{5} + 2\cos\frac{4\pi}{5} = 0 .
        \]
  \item $c = \cos\frac{2\pi}{5}$ रखिए और $\cos\frac{4\pi}{5}$ के द्विकोण सूत्र से प्रश्न 8 को $c$ के लिए
        द्विघात \omterm{def:g10:algebra:equation}{समीकरण} में बदलिए; उसे हल कीजिए और निष्कर्ष निकालिए:
        \[
        \cos 72^\circ = \frac{\sqrt5 - 1}{4} .
        \]
  \item संख्याओं से जाँचिए, फिर स्वर्ण अनुपात को अपने काम पर हस्ताक्षर करने
        दीजिए: दिखाइए कि $c = \frac{1}{2\varphi}$, जहाँ $\varphi = \frac{1 + \sqrt5}{2}$ (\cref{pb:g10:algebra:1}) --- और वह प्रसिद्ध
        ज्यामितीय प्रतिध्वनि बताइए: सम पंचभुज में विकर्ण भुजा का $\varphi$ गुना
        होता है।
\end{enumerate}

\medskip
\noindent\textbf{भाग III --- दे म्वाव्र के छोटे रास्ते।}
\begin{enumerate}[resume]
  \item रचना का फ़ैसला: $\cos 72^\circ$ को केवल वर्गमूल चाहिए (प्रश्न 9), इसलिए पंचभुज
        पटरी और परकार से रचनीय है; जबकि $\cos 20^\circ$ एक अटूट त्रिघात (\cref{pb:g12:trigo:1}) को
        संतुष्ट करता है, इसलिए $60^\circ$ के कोण का त्रिभाजन नहीं हो सकता। उभरती
        हुई कसौटी बताइए (वर्गमूल बनाए जा सकते हैं, घनमूल नहीं), जिसका पूरा
        सिद्धांत --- उन्नीस वर्ष के गाउस और 17-भुज --- स्नातक खंडों में
        रहता है।
  \item $z^4 = -4$ को \omterm{def:g12:complex:argument}{चरघातांकी रूप} से हल कीजिए और \omterm{def:g12:complex:conjugate}{संयुग्मी} \omterm{def:g11:quad:discriminant}{मूलों} की जोड़ी बनाकर
        \cref{exo:g12:complex:8} का \omterm{def:g10:algebra:expand}{गुणनखंडन} $z^4 + 4 = (z^2 - 2z + 2)(z^2 + 2z + 2)$ फिर से पाइए।
  \item द्विपद प्रमेय और दे म्वाव्र से $(\cos\theta + \iu\sin\theta)^3$ का प्रसार कीजिए, और $\cos 3\theta = 4\cos^3\theta - 3\cos\theta$ को
        तीन पंक्तियों में फिर निकालिए --- \cref{pb:g12:trigo:1} के योग-सूत्र वाले रास्ते से
        तुलना कीजिए।
  \item \cref{exo:g12:complex:9} से जाँच: $\theta = \frac{2\pi}{5}$ और $n = 4$ के लिए योग $1 + \cos\theta + \dots + \cos 4\theta$ किसके बराबर है, और वह
        प्रश्न 8 के अनुकूल क्यों है?
  \item $8\iu$ के तीनों घनमूल \omterm{def:g12:complex:def}{बीजगणितीय रूप} में निकालिए।
\end{enumerate}

\medskip
\noindent\textbf{भाग IV --- गुणा ही ज्यामिति है।}
\begin{enumerate}[resume]
  \item प्रतिचित्रण $z \mapsto (1 + \iu)z$ का पूरा वर्णन कीजिए: वह कितने कोण से घुमाता है,
        कितने गुना बड़ा करता है, और इकाई वर्ग के साथ क्या करता है?
  \item संयुग्मियों (\cref{prop:g12:complex:modulus}) से $\abs{z_1 z_2} = \abs{z_1}\,\abs{z_2}$ सिद्ध कीजिए, फिर $n = 1, 2, 3, 4$ के लिए $(1+\iu)^n$
        निकालिए और उन घातों द्वारा खींचे गए सर्पिल का वर्णन कीजिए।
  \item संख्या $\mathrm{j} = \eu^{2\iu\pi/3}$: दिखाइए कि $1 + \mathrm{j} + \mathrm{j}^2 = 0$, और त्रिभुज $\left(1, \mathrm{j}, \mathrm{j}^2\right)$ पर चिरपरिचित कसौटी
        सत्यापित कीजिए: $a + \mathrm{j} b + \mathrm{j}^2 c = 0$ किसी समबाहु त्रिभुज $abc$ (के एक अभिविन्यास)
        के लिए सही होता है।
  \item बीजगणित की मूल प्रमेय (स्वीकृत: दालम्बेर--गाउस): हर बहुपद $\C$ पर
        पूरी तरह गुणनखंडित हो जाता है। इससे वास्तविक संसार वाला उपप्रमेय
        निकालिए --- हर वास्तविक बहुपद वास्तविक रैखिक और द्विघात टुकड़ों में
        गुणनखंडित हो जाता है --- और प्रश्न 12 को उसका उदाहरण बताइए।
  \item समापन --- $\C$ का चित्र, हर बात पर एक वाक्य: संख्याएँ बिंदु बन जाती
        हैं; गुणा घूर्णन-और-मापन बन जाता है; \omterm{def:g12:complex:rootsofunity}{इकाई के मूल} सम बहुभुज बन जाते
        हैं; दे म्वाव्र घातों को त्रिकोणमितीय सर्वसमिकाओं में बदल देते हैं;
        और बीजगणितीय संवृतता यह गारंटी देती है कि किसी \omterm{def:g10:algebra:equation}{समीकरण} को कभी किसी
        बड़े संसार की ज़रूरत नहीं पड़ेगी। समापन-पद: प्रश्न 10 में कौन-सी दो
        प्रसिद्ध संख्याएँ मिलीं?
\end{enumerate}
\end{problem}
